
For a finite 2-vertex-connected graph \(G\), let \(\beta(G)\) be the largest
number of consecutive odd ears at the beginning of an open ear decomposition,
counting its initial circuit and its one-edge ears. We prove that deciding
\(\beta(G)\ge t\), with \(t\) part of the input, is NP-complete even for simple
graphs with a universal vertex whose deletion leaves a connected bipartite
graph. This answers the variable-threshold complexity question posed by
Seb\H{o} and by Benchetrit--Seb\H{o}. The reduction is from explicitly encoded
Set Cover. It uses the classical factor-critical formulation of \(\beta\),
a cone identity, and a block-loss inequality that permits arbitrary partial
blocks. In particular, it does not assume an optimal subgraph is a union of
whole blocks.
